% Generated draft: requires mathematical and source review.
\documentclass[11pt]{article}
\usepackage[margin=25mm]{geometry}
\usepackage{fontspec}
\setmainfont{Latin Modern Roman}
\usepackage{amsmath,amssymb,mathtools,unicode-math}
\setmathfont{Latin Modern Math}
\usepackage{enumitem,needspace,xurl,hyperref}
\setlength{\parindent}{0pt}
\setlength{\parskip}{5pt}
\setlength{\emergencystretch}{3em}
\newcommand{\N}{\mathbb{N}}
\newcommand{\Z}{\mathbb{Z}}
\newcommand{\Q}{\mathbb{Q}}
\newcommand{\R}{\mathbb{R}}
\newcommand{\C}{\mathbb{C}}
\newcommand{\F}{\mathbb{F}}
\newcommand{\A}{\mathbb{A}}
\newcommand{\PP}{\mathbb{P}}
\newcommand{\E}{\mathbb{E}}
\newcommand{\eps}{\varepsilon}
\newcommand{\dd}{\,\mathrm{d}}
\begin{document}
\section*{UnsolvedMath ID: 800251}
\textbf{Source problem code:} catalogue entry
\subsection{Definitions and mathematical statement}
Let \(n\geq 1\) and let \(\Omega\Subset\C^n\) be a pseudoconvex domain with smooth boundary. Write
\[
\mathcal P(\overline\Omega)=\{\phi\in C(\overline\Omega,\R):\phi|_\Omega\text{ is plurisubharmonic}\}.
\]
For a compact set \(K\subseteq\overline\Omega\), define its hull by
\[
\widehat K=\left\{z\in\overline\Omega:\phi(z)\leq\sup_{w\in K}\phi(w)\text{ for every }\phi\in\mathcal P(\overline\Omega)\right\}.
\]
For the empty set, use \(\sup\varnothing=-\infty\), so that \(\widehat\varnothing=\varnothing\).
A Stein neighborhood basis of \(\overline\Omega\) means that, for every open set \(V\subseteq\C^n\) containing \(\overline\Omega\), there exists an open Stein set \(U\subseteq\C^n\) such that
\[
\overline\Omega\subseteq U\subseteq V.
\]
The natural universally quantified reading of the supplied assertion is the following conjectural implication:
\[
\left[\text{for every compact }K\subseteq\overline\Omega,\quad
\widehat K\cap\partial\Omega=K\cap\partial\Omega\right]
\quad\Longrightarrow\quad
\overline\Omega\text{ has a Stein neighborhood basis}.
\]
The source introduces a compact set \(K\) and then assumes the displayed boundary equality, without explicitly stating that the equality must hold for every compact \(K\). The universal quantifier in this formulation therefore requires human confirmation. An assumption concerning only a single chosen compact set would be a different assertion. The source also leaves the precise differentiability convention for ``smooth boundary'' unstated.
\subsection{Short English statement}
Under the universal reading of the hypothesis, the question asks whether a relatively compact smoothly bounded pseudoconvex domain has a Stein neighborhood basis whenever the hull of every compact subset of its closure acquires no additional boundary points. The quantifier over compact subsets requires confirmation from the source. The supplied registry labels the problem open; this status has not been independently verified.
\subsection{Sources}
\begin{itemize}
\item Local UnsolvedMath registry, permanent ID \(800251\), source problem code \texttt{catalogue entry}, titled ``Several Complex Variables Problem List --- Problem 50.'' The supplied record labels the problem \texttt{open}.
\item Buzzard, Fornaess, and Krantz, \emph{The Several Complex Variables Problem List} (1995), Problem 50. Supplied source URL: \url{https://arxiv.org/abs/math/9509202}. The record identifies the extraction method as \texttt{arxiv\_latex} and records access on 2026-10-16; these are supplied provenance details, not an independently verified visit.
\item The extracted problem is attributed to N.~Sibony, with the supplied contact address \texttt{sibony@anh.matups.fr}. The registry's \texttt{proposed\_by} field lists Buzzard, Fornaess, and Krantz, and its \texttt{proposed\_year} field is 1995; the distinction between list attribution and the attribution attached to the individual problem is retained.
\end{itemize}
\end{document}
